
\documentclass{article}
\usepackage{colortbl}
\usepackage{makecell}
\usepackage{multirow}
\usepackage{supertabular}

\begin{document}

\newcounter{utterance}

\centering \large Interaction Transcript for game `dond', experiment `semi\_en', episode 13 with reka{-}flash{-}3.1{-}t0.0.
\vspace{24pt}

{ \footnotesize  \setcounter{utterance}{1}
\setlength{\tabcolsep}{0pt}
\begin{supertabular}{c@{$\;$}|p{.15\linewidth}@{}p{.15\linewidth}p{.15\linewidth}p{.15\linewidth}p{.15\linewidth}p{.15\linewidth}}
    \# & $\;$A & \multicolumn{4}{c}{Game Master} & $\;\:$B\\
    \hline

    \theutterance \stepcounter{utterance}  
    & & \multicolumn{4}{p{0.6\linewidth}}{
        \cellcolor[rgb]{0.9,0.9,0.9}{
            \makecell[{{p{\linewidth}}}]{
                \texttt{\tiny{[P1$\langle$GM]}}
                \texttt{You are playing a negotiation game in which you have to agree on how to divide a set of items among you and another player.} \\
\\ 
\texttt{Rules:} \\
\texttt{(a) You and the other player are given a set of items. Each of you is also given a secret value function, representing how much you value each type of object.} \\
\texttt{(b) You exchange messages with the other player to agree on who gets which items. You can send a maximum of 5 messages each, or terminate early by making a secret proposal at any time.} \\
\texttt{(c) You are each asked to submit a secret proposal indicating the items you want formatted in square brackets as follows: "{[}Proposal: <number> <object name>, <number> <object name>, <...>{]}"} \\
\texttt{(d) If your proposals are complementary, i.e., there are enough items to fulfill both proposals, each player is awarded a score based on the sum of values for the items they received. Otherwise, both of you get zero points.} \\
\texttt{(e) Your goal is to maximize the score you receive.} \\
\\ 
\texttt{Let us start.} \\
\\ 
\texttt{The set of available items is:} \\
\\ 
\texttt{4 papers, 2 clocks, 2 potatoes} \\
\\ 
\texttt{Your secret value function is:} \\
\\ 
\texttt{paper: 1, clock: 1, potato: 2} \\
\\ 
\texttt{IMPORTANT: Your messages, unless it is the secret proposal, are directly transmitted to the other player, so do not include any response to the rules or text announcing your message. To make a secret proposal, use the indicated format. Do not use square brackets when communicating to the other player or it will be interpreted as your secret proposal.} \\
            }
        }
    }
    & \\ \\

    \theutterance \stepcounter{utterance}  
    & \multicolumn{4}{p{0.6\linewidth}}{
        \cellcolor[rgb]{0.9,0.9,0.9}{
            \makecell[{{p{\linewidth}}}]{
                \texttt{\tiny{[P1$\rangle$GM]}}
                \texttt{<reasoning>} \\
\texttt{Let's see. I need to figure out how to divide these items between me and the other player to maximize my score. The items are 4 papers, 2 clocks, and 2 potatoes. My values are paper: 1, clock: 1, potato: 2. The other player has their own secret values, which I don't know. The goal is to make a proposal that gives me the highest possible score, considering that the other player will also make a proposal, and if our proposals are complementary (i.e., together they take all the items without overlap), then we each get the sum of the values of what we got. Otherwise, both get zero.} \\
\\ 
\texttt{First, I need to understand what complementary proposals mean. If I propose some items, the other player's proposal must include the remaining items. So, for example, if I take 2 papers and 1 potato, the other player must take the remaining 2 papers, 1 clock, and 1 potato. But if there's any overlap (like both proposing a paper), then it's invalid, and both get zero.} \\
\\ 
\texttt{So my strategy should be to propose a combination of items that I value highly, but also leave enough items for the other player so that their proposal can complement mine. But since I don't know their values, I have to make an assumption. Maybe I should aim for a split that gives me the maximum possible value, assuming the other player will take the least valuable items they can. Or maybe the other player is trying to maximize their own value, so they might take the most valuable items they can. Hmm.} \\
\\ 
\texttt{Wait, but since both are trying to maximize their own scores, perhaps the best approach is to propose a split where I get the items I value the most, and leave the rest. But how to balance that with what the other player might take.} \\
\\ 
\texttt{Let me calculate my total possible value. If I take all items, my value would be 4*1 + 2*1 + 2*2 = 4 + 2 + 4 = 10. But obviously, the other player won't let me take everything. So I need to find a split where I get a good portion.} \\
\\ 
\texttt{Alternatively, maybe the optimal split is where each player gets a bundle that's worth at least half of the total possible value. But since the other player's values are unknown, this is tricky.} \\
\\ 
\texttt{Alternatively, think in terms of} \\
            }
        }
    }
    & & \\ \\

    \theutterance \stepcounter{utterance}  
    & & & \multicolumn{2}{p{0.3\linewidth}}{
        \cellcolor[rgb]{0.9,0.9,0.9}{
            \makecell[{{p{\linewidth}}}]{
                \texttt{\tiny{[GM$|$GM]}}
                \texttt{continue} \\
            }
        }
    }
    & & \\ \\

    \theutterance \stepcounter{utterance}  
    & & & \multicolumn{2}{p{0.3\linewidth}}{
        \cellcolor[rgb]{0.9,0.9,0.9}{
            \makecell[{{p{\linewidth}}}]{
                \texttt{\tiny{[GM$|$GM]}}
                \texttt{<reasoning>} \\
\texttt{Let's see. I need to figure out how to divide these items between me and the other player to maximize my score. The items are 4 papers, 2 clocks, and 2 potatoes. My values are paper: 1, clock: 1, potato: 2. The other player has their own secret values, which I don't know. The goal is to make a proposal that gives me the highest possible score, considering that the other player will also make a proposal, and if our proposals are complementary (i.e., together they take all the items without overlap), then we each get the sum of the values of what we got. Otherwise, both get zero.} \\
\\ 
\texttt{First, I need to understand what complementary proposals mean. If I propose some items, the other player's proposal must include the remaining items. So, for example, if I take 2 papers and 1 potato, the other player must take the remaining 2 papers, 1 clock, and 1 potato. But if there's any overlap (like both proposing a paper), then it's invalid, and both get zero.} \\
\\ 
\texttt{So my strategy should be to propose a combination of items that I value highly, but also leave enough items for the other player so that their proposal can complement mine. But since I don't know their values, I have to make an assumption. Maybe I should aim for a split that gives me the maximum possible value, assuming the other player will take the least valuable items they can. Or maybe the other player is trying to maximize their own value, so they might take the most valuable items they can. Hmm.} \\
\\ 
\texttt{Wait, but since both are trying to maximize their own scores, perhaps the best approach is to propose a split where I get the items I value the most, and leave the rest. But how to balance that with what the other player might take.} \\
\\ 
\texttt{Let me calculate my total possible value. If I take all items, my value would be 4*1 + 2*1 + 2*2 = 4 + 2 + 4 = 10. But obviously, the other player won't let me take everything. So I need to find a split where I get a good portion.} \\
\\ 
\texttt{Alternatively, maybe the optimal split is where each player gets a bundle that's worth at least half of the total possible value. But since the other player's values are unknown, this is tricky.} \\
\\ 
\texttt{Alternatively, think in terms of} \\
            }
        }
    }
    & & \\ \\

    \theutterance \stepcounter{utterance}  
    & & \multicolumn{4}{p{0.6\linewidth}}{
        \cellcolor[rgb]{0.9,0.9,0.9}{
            \makecell[{{p{\linewidth}}}]{
                \texttt{\tiny{[GM$\rangle$P2]}}
                \texttt{You are playing a negotiation game in which you have to agree on how to divide a set of items among you and another player.} \\
\\ 
\texttt{Rules:} \\
\texttt{(a) You and the other player are given a set of items. Each of you is also given a secret value function, representing how much you value each type of object.} \\
\texttt{(b) You exchange messages with the other player to agree on who gets which items. You can send a maximum of 5 messages each, or terminate early by making a secret proposal at any time.} \\
\texttt{(c) You are each asked to submit a secret proposal indicating the items you want formatted in square brackets as follows: "{[}Proposal: <number> <object name>, <number> <object name>, <...>{]}"} \\
\texttt{(d) If your proposals are complementary, i.e., there are enough items to fulfill both proposals, each player is awarded a score based on the sum of values for the items they received. Otherwise, both of you get zero points.} \\
\texttt{(e) Your goal is to maximize the score you receive.} \\
\\ 
\texttt{Let us start.} \\
\\ 
\texttt{The set of available items is:} \\
\\ 
\texttt{4 papers, 2 clocks, 2 potatoes} \\
\\ 
\texttt{Your secret value function is:} \\
\\ 
\texttt{paper: 0, clock: 4, potato: 1} \\
\\ 
\texttt{IMPORTANT: Your messages, unless it is the secret proposal, are directly transmitted to the other player, so do not include any response to the rules or text announcing your message. To make a secret proposal, use the indicated format. Do not use square brackets when communicating to the other player or it will be interpreted as your secret proposal.} \\
\\ 
\\ 
\texttt{<reasoning>} \\
\texttt{Let's see. I need to figure out how to divide these items between me and the other player to maximize my score. The items are 4 papers, 2 clocks, and 2 potatoes. My values are paper: 1, clock: 1, potato: 2. The other player has their own secret values, which I don't know. The goal is to make a proposal that gives me the highest possible score, considering that the other player will also make a proposal, and if our proposals are complementary (i.e., together they take all the items without overlap), then we each get the sum of the values of what we got. Otherwise, both get zero.} \\
\\ 
\texttt{First, I need to understand what complementary proposals mean. If I propose some items, the other player's proposal must include the remaining items. So, for example, if I take 2 papers and 1 potato, the other player must take the remaining 2 papers, 1 clock, and 1 potato. But if there's any overlap (like both proposing a paper), then it's invalid, and both get zero.} \\
\\ 
\texttt{So my strategy should be to propose a combination of items that I value highly, but also leave enough items for the other player so that their proposal can complement mine. But since I don't know their values, I have to make an assumption. Maybe I should aim for a split that gives me the maximum possible value, assuming the other player will take the least valuable items they can. Or maybe the other player is trying to maximize their own value, so they might take the most valuable items they can. Hmm.} \\
\\ 
\texttt{Wait, but since both are trying to maximize their own scores, perhaps the best approach is to propose a split where I get the items I value the most, and leave the rest. But how to balance that with what the other player might take.} \\
\\ 
\texttt{Let me calculate my total possible value. If I take all items, my value would be 4*1 + 2*1 + 2*2 = 4 + 2 + 4 = 10. But obviously, the other player won't let me take everything. So I need to find a split where I get a good portion.} \\
\\ 
\texttt{Alternatively, maybe the optimal split is where each player gets a bundle that's worth at least half of the total possible value. But since the other player's values are unknown, this is tricky.} \\
\\ 
\texttt{Alternatively, think in terms of} \\
            }
        }
    }
    & \\ \\

    \theutterance \stepcounter{utterance}  
    & & & \multicolumn{4}{p{0.6\linewidth}}{
        \cellcolor[rgb]{0.9,0.9,0.9}{
            \makecell[{{p{\linewidth}}}]{
                \texttt{\tiny{[GM$\langle$P2]}}
                \texttt{<reasoning>} \\
\texttt{Let me try to figure out the best way to approach this negotiation. The items available are 4 papers, 2 clocks, and 2 potatoes. My value for each is paper: 1, clock: 1, potato: 2. The other player has their own values, which I don't know. The key is to make a proposal that gives me the highest possible score, considering that the other player will also propose, and if our proposals are complementary, we each get the sum of the values of our items. If not, both get zero.} \\
\\ 
\texttt{First, I need to understand what complementary means. If I propose some items, the other player's proposal must include exactly the remaining items. So, for example, if I take 2 papers and 1 potato, the other player must take the remaining 2 papers, 1 clock, and 1 potato. But if there's any overlap (like both proposing a paper), it's invalid.} \\
\\ 
\texttt{So my strategy should be to propose a combination of items that I value highly, but also leave enough items for the other player so that their proposal can complement mine. But since I don't know their values, I have to make an assumption. Maybe the other player will try to maximize their own value, so they might take the most valuable items they can. But since I don't know their values, perhaps I should aim for a split that gives me the maximum possible value, assuming the other player will take the least valuable items they can. Or maybe the other player is trying to maximize their own value, so they might take the most valuable items they can. Hmm.} \\
\\ 
\texttt{Let me calculate my total possible value. If I take all items, my value would be 4*1 + 2*1 + 2*2 = 4 + 2 + 4 = 10. But obviously, the other player won't let me take everything. So I need to find a split where I get a good portion.} \\
\\ 
\texttt{Alternatively, maybe the optimal split is where each player gets a bundle that's worth at least half of the total possible value. But since the other player's values are unknown, this is tricky.} \\
\\ 
\texttt{Alternatively, think in terms of what items are most valuable to me. Potatoes are worth 2 each, clocks 1 each, papers 1 each. So I should prioritize getting potatoes. There are 2 potatoes. If I take both potatoes, that's} \\
            }
        }
    }
     \\ \\

    \theutterance \stepcounter{utterance}  
    & & & \multicolumn{2}{p{0.3\linewidth}}{
        \cellcolor[rgb]{0.9,0.9,0.9}{
            \makecell[{{p{\linewidth}}}]{
                \texttt{\tiny{[GM$|$GM]}}
                \texttt{continue} \\
            }
        }
    }
    & & \\ \\

    \theutterance \stepcounter{utterance}  
    & & & \multicolumn{2}{p{0.3\linewidth}}{
        \cellcolor[rgb]{0.9,0.9,0.9}{
            \makecell[{{p{\linewidth}}}]{
                \texttt{\tiny{[GM$|$GM]}}
                \texttt{<reasoning>} \\
\texttt{Let me try to figure out the best way to approach this negotiation. The items available are 4 papers, 2 clocks, and 2 potatoes. My value for each is paper: 1, clock: 1, potato: 2. The other player has their own values, which I don't know. The key is to make a proposal that gives me the highest possible score, considering that the other player will also propose, and if our proposals are complementary, we each get the sum of the values of our items. If not, both get zero.} \\
\\ 
\texttt{First, I need to understand what complementary means. If I propose some items, the other player's proposal must include exactly the remaining items. So, for example, if I take 2 papers and 1 potato, the other player must take the remaining 2 papers, 1 clock, and 1 potato. But if there's any overlap (like both proposing a paper), it's invalid.} \\
\\ 
\texttt{So my strategy should be to propose a combination of items that I value highly, but also leave enough items for the other player so that their proposal can complement mine. But since I don't know their values, I have to make an assumption. Maybe the other player will try to maximize their own value, so they might take the most valuable items they can. But since I don't know their values, perhaps I should aim for a split that gives me the maximum possible value, assuming the other player will take the least valuable items they can. Or maybe the other player is trying to maximize their own value, so they might take the most valuable items they can. Hmm.} \\
\\ 
\texttt{Let me calculate my total possible value. If I take all items, my value would be 4*1 + 2*1 + 2*2 = 4 + 2 + 4 = 10. But obviously, the other player won't let me take everything. So I need to find a split where I get a good portion.} \\
\\ 
\texttt{Alternatively, maybe the optimal split is where each player gets a bundle that's worth at least half of the total possible value. But since the other player's values are unknown, this is tricky.} \\
\\ 
\texttt{Alternatively, think in terms of what items are most valuable to me. Potatoes are worth 2 each, clocks 1 each, papers 1 each. So I should prioritize getting potatoes. There are 2 potatoes. If I take both potatoes, that's} \\
            }
        }
    }
    & & \\ \\

    \theutterance \stepcounter{utterance}  
    & & \multicolumn{4}{p{0.6\linewidth}}{
        \cellcolor[rgb]{0.9,0.9,0.9}{
            \makecell[{{p{\linewidth}}}]{
                \texttt{\tiny{[P1$\langle$GM]}}
                \texttt{<reasoning>} \\
\texttt{Let me try to figure out the best way to approach this negotiation. The items available are 4 papers, 2 clocks, and 2 potatoes. My value for each is paper: 1, clock: 1, potato: 2. The other player has their own values, which I don't know. The key is to make a proposal that gives me the highest possible score, considering that the other player will also propose, and if our proposals are complementary, we each get the sum of the values of our items. If not, both get zero.} \\
\\ 
\texttt{First, I need to understand what complementary means. If I propose some items, the other player's proposal must include exactly the remaining items. So, for example, if I take 2 papers and 1 potato, the other player must take the remaining 2 papers, 1 clock, and 1 potato. But if there's any overlap (like both proposing a paper), it's invalid.} \\
\\ 
\texttt{So my strategy should be to propose a combination of items that I value highly, but also leave enough items for the other player so that their proposal can complement mine. But since I don't know their values, I have to make an assumption. Maybe the other player will try to maximize their own value, so they might take the most valuable items they can. But since I don't know their values, perhaps I should aim for a split that gives me the maximum possible value, assuming the other player will take the least valuable items they can. Or maybe the other player is trying to maximize their own value, so they might take the most valuable items they can. Hmm.} \\
\\ 
\texttt{Let me calculate my total possible value. If I take all items, my value would be 4*1 + 2*1 + 2*2 = 4 + 2 + 4 = 10. But obviously, the other player won't let me take everything. So I need to find a split where I get a good portion.} \\
\\ 
\texttt{Alternatively, maybe the optimal split is where each player gets a bundle that's worth at least half of the total possible value. But since the other player's values are unknown, this is tricky.} \\
\\ 
\texttt{Alternatively, think in terms of what items are most valuable to me. Potatoes are worth 2 each, clocks 1 each, papers 1 each. So I should prioritize getting potatoes. There are 2 potatoes. If I take both potatoes, that's} \\
            }
        }
    }
    & \\ \\

    \theutterance \stepcounter{utterance}  
    & \multicolumn{4}{p{0.6\linewidth}}{
        \cellcolor[rgb]{0.9,0.9,0.9}{
            \makecell[{{p{\linewidth}}}]{
                \texttt{\tiny{[P1$\rangle$GM]}}
                \texttt{<reasoning>} \\
\texttt{Let me try to figure out the best way to divide these items. The items are 4 papers, 2 clocks, and 2 potatoes. My values are paper:1, clock:1, potato:2. The other player has their own values, which I don't know. The goal is to make a proposal that gives me the highest possible score if our proposals are complementary.} \\
\\ 
\texttt{First, I need to maximize my own value. Potatoes are the most valuable to me, so I should aim to get as many as possible. There are 2 potatoes. If I take both, that's 2*2=4 points. Then, maybe take some papers or clocks. But the other player might also want potatoes. Since there are only 2, if I take both, the other player can't get any. But maybe the other player values potatoes more than I do? I don't know. But since I know my own values, maybe I should take the potatoes and see.} \\
\\ 
\texttt{Alternatively, maybe split the potatoes. If I take 1 potato, the other player takes 1. Then I can get more other items. Let's see. If I take 1 potato (2 points), then maybe take 2 papers (2 points) and 1 clock (1 point). Total 5. But the other player would get the remaining 1 potato (2), 2 papers, 1 clock. Their total would be 2 + 2 +1=5. But if their values are different, maybe they can get more. But since I don't know, maybe equal splits are safer? Not sure.} \\
\\ 
\texttt{Wait, but the other player is also trying to maximize their own score. So if I take both potatoes, their maximum possible is 4 papers + 2 clocks = 4+2=6. But if I take both potatoes (4 points), and maybe some other items, then my total would be higher. Let's see: If I take both potatoes (4), 2 papers (2), and 1 clock (1). Total 7. The other player would get 2 papers and 1 clock. Their value depends on their own values. If they value papers and clocks more than potatoes, maybe they can get more. But if they value potatoes more, they might get less. But since I don't know, maybe taking both potatoes is risky because the other player might not get anything if} \\
            }
        }
    }
    & & \\ \\

    \theutterance \stepcounter{utterance}  
    & & & \multicolumn{2}{p{0.3\linewidth}}{
        \cellcolor[rgb]{0.9,0.9,0.9}{
            \makecell[{{p{\linewidth}}}]{
                \texttt{\tiny{[GM$|$GM]}}
                \texttt{continue} \\
            }
        }
    }
    & & \\ \\

    \theutterance \stepcounter{utterance}  
    & & & \multicolumn{2}{p{0.3\linewidth}}{
        \cellcolor[rgb]{0.9,0.9,0.9}{
            \makecell[{{p{\linewidth}}}]{
                \texttt{\tiny{[GM$|$GM]}}
                \texttt{<reasoning>} \\
\texttt{Let me try to figure out the best way to divide these items. The items are 4 papers, 2 clocks, and 2 potatoes. My values are paper:1, clock:1, potato:2. The other player has their own values, which I don't know. The goal is to make a proposal that gives me the highest possible score if our proposals are complementary.} \\
\\ 
\texttt{First, I need to maximize my own value. Potatoes are the most valuable to me, so I should aim to get as many as possible. There are 2 potatoes. If I take both, that's 2*2=4 points. Then, maybe take some papers or clocks. But the other player might also want potatoes. Since there are only 2, if I take both, the other player can't get any. But maybe the other player values potatoes more than I do? I don't know. But since I know my own values, maybe I should take the potatoes and see.} \\
\\ 
\texttt{Alternatively, maybe split the potatoes. If I take 1 potato, the other player takes 1. Then I can get more other items. Let's see. If I take 1 potato (2 points), then maybe take 2 papers (2 points) and 1 clock (1 point). Total 5. But the other player would get the remaining 1 potato (2), 2 papers, 1 clock. Their total would be 2 + 2 +1=5. But if their values are different, maybe they can get more. But since I don't know, maybe equal splits are safer? Not sure.} \\
\\ 
\texttt{Wait, but the other player is also trying to maximize their own score. So if I take both potatoes, their maximum possible is 4 papers + 2 clocks = 4+2=6. But if I take both potatoes (4 points), and maybe some other items, then my total would be higher. Let's see: If I take both potatoes (4), 2 papers (2), and 1 clock (1). Total 7. The other player would get 2 papers and 1 clock. Their value depends on their own values. If they value papers and clocks more than potatoes, maybe they can get more. But if they value potatoes more, they might get less. But since I don't know, maybe taking both potatoes is risky because the other player might not get anything if} \\
            }
        }
    }
    & & \\ \\

    \theutterance \stepcounter{utterance}  
    & & \multicolumn{4}{p{0.6\linewidth}}{
        \cellcolor[rgb]{0.9,0.9,0.9}{
            \makecell[{{p{\linewidth}}}]{
                \texttt{\tiny{[GM$\rangle$P2]}}
                \texttt{<reasoning>} \\
\texttt{Let me try to figure out the best way to divide these items. The items are 4 papers, 2 clocks, and 2 potatoes. My values are paper:1, clock:1, potato:2. The other player has their own values, which I don't know. The goal is to make a proposal that gives me the highest possible score if our proposals are complementary.} \\
\\ 
\texttt{First, I need to maximize my own value. Potatoes are the most valuable to me, so I should aim to get as many as possible. There are 2 potatoes. If I take both, that's 2*2=4 points. Then, maybe take some papers or clocks. But the other player might also want potatoes. Since there are only 2, if I take both, the other player can't get any. But maybe the other player values potatoes more than I do? I don't know. But since I know my own values, maybe I should take the potatoes and see.} \\
\\ 
\texttt{Alternatively, maybe split the potatoes. If I take 1 potato, the other player takes 1. Then I can get more other items. Let's see. If I take 1 potato (2 points), then maybe take 2 papers (2 points) and 1 clock (1 point). Total 5. But the other player would get the remaining 1 potato (2), 2 papers, 1 clock. Their total would be 2 + 2 +1=5. But if their values are different, maybe they can get more. But since I don't know, maybe equal splits are safer? Not sure.} \\
\\ 
\texttt{Wait, but the other player is also trying to maximize their own score. So if I take both potatoes, their maximum possible is 4 papers + 2 clocks = 4+2=6. But if I take both potatoes (4 points), and maybe some other items, then my total would be higher. Let's see: If I take both potatoes (4), 2 papers (2), and 1 clock (1). Total 7. The other player would get 2 papers and 1 clock. Their value depends on their own values. If they value papers and clocks more than potatoes, maybe they can get more. But if they value potatoes more, they might get less. But since I don't know, maybe taking both potatoes is risky because the other player might not get anything if} \\
            }
        }
    }
    & \\ \\

    \theutterance \stepcounter{utterance}  
    & & & \multicolumn{4}{p{0.6\linewidth}}{
        \cellcolor[rgb]{0.9,0.9,0.9}{
            \makecell[{{p{\linewidth}}}]{
                \texttt{\tiny{[GM$\langle$P2]}}
                \texttt{<reasoning>} \\
\texttt{Let me try to work through this step by step. The goal is to maximize my score by proposing a split of the items such that my proposal and the other player's are complementary. Since I don't know their values, I need to make an educated guess.} \\
\\ 
\texttt{First, let's list out the items and my values:} \\
\\ 
\texttt{{-} 4 papers, each worth 1 to me.} \\
\texttt{{-} 2 clocks, each worth 1 to me.} \\
\texttt{{-} 2 potatoes, each worth 2 to me.} \\
\\ 
\texttt{Total value if I take everything: 4*1 + 2*1 + 2*2 = 4 + 2 + 4 = 10. But obviously, the other player won't let me have everything. So I need to find a balance.} \\
\\ 
\texttt{Potatoes are the most valuable to me. There are only 2, so if I take both, that's 4 points. Then I can take some papers or clocks. Let's see: If I take both potatoes (4), 2 papers (2), and 1 clock (1), that's 7 points. The other player would get the remaining 2 papers and 1 clock. Their value depends on their preferences, but if they value papers and clocks higher than potatoes, they might get more. But if they value potatoes, they might get less. Since I don't know, maybe taking both potatoes is a good start.} \\
\\ 
\texttt{Alternatively, if I take 1 potato (2), then maybe take more papers or clocks. Let's say 1 potato, 3 papers, 1 clock. That's 2 + 3 +1 =6. The other player gets 1 potato, 1 paper, 1 clock. Their value could be higher if they value potatoes more. But I'm not sure.} \\
\\ 
\texttt{Wait, but the other player is also trying to maximize their own score. So if I take both potatoes, they might not get any, which could lead to their proposal not complementing mine if they also want potatoes. But if I leave one potato, maybe they take it, and we can split the rest.} \\
\\ 
\texttt{Alternatively, maybe the best approach is to propose a split where I get the most valuable items I can, and leave the rest. Since potatoes are highest, take both. Then take as many papers or clocks as possible. Let's see:} \\
\\ 
\texttt{If I take both potatoes (2 each =4), 2 papers (1 each=2} \\
            }
        }
    }
     \\ \\

    \theutterance \stepcounter{utterance}  
    & & & \multicolumn{2}{p{0.3\linewidth}}{
        \cellcolor[rgb]{0.9,0.9,0.9}{
            \makecell[{{p{\linewidth}}}]{
                \texttt{\tiny{[GM$|$GM]}}
                \texttt{continue} \\
            }
        }
    }
    & & \\ \\

    \theutterance \stepcounter{utterance}  
    & & & \multicolumn{2}{p{0.3\linewidth}}{
        \cellcolor[rgb]{0.9,0.9,0.9}{
            \makecell[{{p{\linewidth}}}]{
                \texttt{\tiny{[GM$|$GM]}}
                \texttt{<reasoning>} \\
\texttt{Let me try to work through this step by step. The goal is to maximize my score by proposing a split of the items such that my proposal and the other player's are complementary. Since I don't know their values, I need to make an educated guess.} \\
\\ 
\texttt{First, let's list out the items and my values:} \\
\\ 
\texttt{{-} 4 papers, each worth 1 to me.} \\
\texttt{{-} 2 clocks, each worth 1 to me.} \\
\texttt{{-} 2 potatoes, each worth 2 to me.} \\
\\ 
\texttt{Total value if I take everything: 4*1 + 2*1 + 2*2 = 4 + 2 + 4 = 10. But obviously, the other player won't let me have everything. So I need to find a balance.} \\
\\ 
\texttt{Potatoes are the most valuable to me. There are only 2, so if I take both, that's 4 points. Then I can take some papers or clocks. Let's see: If I take both potatoes (4), 2 papers (2), and 1 clock (1), that's 7 points. The other player would get the remaining 2 papers and 1 clock. Their value depends on their preferences, but if they value papers and clocks higher than potatoes, they might get more. But if they value potatoes, they might get less. Since I don't know, maybe taking both potatoes is a good start.} \\
\\ 
\texttt{Alternatively, if I take 1 potato (2), then maybe take more papers or clocks. Let's say 1 potato, 3 papers, 1 clock. That's 2 + 3 +1 =6. The other player gets 1 potato, 1 paper, 1 clock. Their value could be higher if they value potatoes more. But I'm not sure.} \\
\\ 
\texttt{Wait, but the other player is also trying to maximize their own score. So if I take both potatoes, they might not get any, which could lead to their proposal not complementing mine if they also want potatoes. But if I leave one potato, maybe they take it, and we can split the rest.} \\
\\ 
\texttt{Alternatively, maybe the best approach is to propose a split where I get the most valuable items I can, and leave the rest. Since potatoes are highest, take both. Then take as many papers or clocks as possible. Let's see:} \\
\\ 
\texttt{If I take both potatoes (2 each =4), 2 papers (1 each=2} \\
            }
        }
    }
    & & \\ \\

    \theutterance \stepcounter{utterance}  
    & & \multicolumn{4}{p{0.6\linewidth}}{
        \cellcolor[rgb]{0.9,0.9,0.9}{
            \makecell[{{p{\linewidth}}}]{
                \texttt{\tiny{[P1$\langle$GM]}}
                \texttt{<reasoning>} \\
\texttt{Let me try to work through this step by step. The goal is to maximize my score by proposing a split of the items such that my proposal and the other player's are complementary. Since I don't know their values, I need to make an educated guess.} \\
\\ 
\texttt{First, let's list out the items and my values:} \\
\\ 
\texttt{{-} 4 papers, each worth 1 to me.} \\
\texttt{{-} 2 clocks, each worth 1 to me.} \\
\texttt{{-} 2 potatoes, each worth 2 to me.} \\
\\ 
\texttt{Total value if I take everything: 4*1 + 2*1 + 2*2 = 4 + 2 + 4 = 10. But obviously, the other player won't let me have everything. So I need to find a balance.} \\
\\ 
\texttt{Potatoes are the most valuable to me. There are only 2, so if I take both, that's 4 points. Then I can take some papers or clocks. Let's see: If I take both potatoes (4), 2 papers (2), and 1 clock (1), that's 7 points. The other player would get the remaining 2 papers and 1 clock. Their value depends on their preferences, but if they value papers and clocks higher than potatoes, they might get more. But if they value potatoes, they might get less. Since I don't know, maybe taking both potatoes is a good start.} \\
\\ 
\texttt{Alternatively, if I take 1 potato (2), then maybe take more papers or clocks. Let's say 1 potato, 3 papers, 1 clock. That's 2 + 3 +1 =6. The other player gets 1 potato, 1 paper, 1 clock. Their value could be higher if they value potatoes more. But I'm not sure.} \\
\\ 
\texttt{Wait, but the other player is also trying to maximize their own score. So if I take both potatoes, they might not get any, which could lead to their proposal not complementing mine if they also want potatoes. But if I leave one potato, maybe they take it, and we can split the rest.} \\
\\ 
\texttt{Alternatively, maybe the best approach is to propose a split where I get the most valuable items I can, and leave the rest. Since potatoes are highest, take both. Then take as many papers or clocks as possible. Let's see:} \\
\\ 
\texttt{If I take both potatoes (2 each =4), 2 papers (1 each=2} \\
            }
        }
    }
    & \\ \\

    \theutterance \stepcounter{utterance}  
    & \multicolumn{4}{p{0.6\linewidth}}{
        \cellcolor[rgb]{0.9,0.9,0.9}{
            \makecell[{{p{\linewidth}}}]{
                \texttt{\tiny{[P1$\rangle$GM]}}
                \texttt{<reasoning>} \\
\texttt{Let me try to figure out the best way to divide these items. The items are 4 papers, 2 clocks, and 2 potatoes. My values are paper:1, clock:1, potato:2. The other player has their own values, which I don't know. The goal is to make a proposal that gives me the highest possible score if our proposals are complementary.} \\
\\ 
\texttt{First, I need to maximize my own value. Potatoes are the most valuable to me, so I should aim to get as many as possible. There are 2 potatoes. If I take both, that's 2*2=4 points. Then, maybe take some papers or clocks. Let's see: If I take both potatoes (4), 2 papers (2), and 1 clock (1). Total 7. The other player would get 2 papers and 1 clock. Their value depends on their own values. If they value papers and clocks more than potatoes, maybe they can get more. But if they value potatoes more, they might get less. But since I don't know, maybe taking both potatoes is risky because the other player might not get anything if they also want potatoes. But if I leave one potato, maybe they take it, and we can split the rest.} \\
\\ 
\texttt{Alternatively, if I take 1 potato (2), then maybe take more papers or clocks. Let's say 1 potato, 3 papers, 1 clock. That's 2 + 3 +1 =6. The other player gets 1 potato, 1 paper, 1 clock. Their value could be higher if they value potatoes more. But I'm not sure.} \\
\\ 
\texttt{Wait, but the other player is also trying to maximize their own score. So if I take both potatoes, they might not get any, which could lead to their proposal not complementing mine if they also want potatoes. But if I leave one potato, maybe they take it, and we can split the rest.} \\
\\ 
\texttt{Alternatively, maybe the best approach is to propose a split where I get the most valuable items I can, and leave the rest. Since potatoes are highest, take both. Then take as many papers or clocks as possible. Let's see:} \\
\\ 
\texttt{If I take both potatoes (2 each =4), 2 papers (1 each=2), and 1 clock (1). Total 7. The other player would get 2 papers and 1 clock} \\
            }
        }
    }
    & & \\ \\

    \theutterance \stepcounter{utterance}  
    & & & \multicolumn{2}{p{0.3\linewidth}}{
        \cellcolor[rgb]{0.9,0.9,0.9}{
            \makecell[{{p{\linewidth}}}]{
                \texttt{\tiny{[GM$|$GM]}}
                \texttt{continue} \\
            }
        }
    }
    & & \\ \\

    \theutterance \stepcounter{utterance}  
    & & & \multicolumn{2}{p{0.3\linewidth}}{
        \cellcolor[rgb]{0.9,0.9,0.9}{
            \makecell[{{p{\linewidth}}}]{
                \texttt{\tiny{[GM$|$GM]}}
                \texttt{<reasoning>} \\
\texttt{Let me try to figure out the best way to divide these items. The items are 4 papers, 2 clocks, and 2 potatoes. My values are paper:1, clock:1, potato:2. The other player has their own values, which I don't know. The goal is to make a proposal that gives me the highest possible score if our proposals are complementary.} \\
\\ 
\texttt{First, I need to maximize my own value. Potatoes are the most valuable to me, so I should aim to get as many as possible. There are 2 potatoes. If I take both, that's 2*2=4 points. Then, maybe take some papers or clocks. Let's see: If I take both potatoes (4), 2 papers (2), and 1 clock (1). Total 7. The other player would get 2 papers and 1 clock. Their value depends on their own values. If they value papers and clocks more than potatoes, maybe they can get more. But if they value potatoes more, they might get less. But since I don't know, maybe taking both potatoes is risky because the other player might not get anything if they also want potatoes. But if I leave one potato, maybe they take it, and we can split the rest.} \\
\\ 
\texttt{Alternatively, if I take 1 potato (2), then maybe take more papers or clocks. Let's say 1 potato, 3 papers, 1 clock. That's 2 + 3 +1 =6. The other player gets 1 potato, 1 paper, 1 clock. Their value could be higher if they value potatoes more. But I'm not sure.} \\
\\ 
\texttt{Wait, but the other player is also trying to maximize their own score. So if I take both potatoes, they might not get any, which could lead to their proposal not complementing mine if they also want potatoes. But if I leave one potato, maybe they take it, and we can split the rest.} \\
\\ 
\texttt{Alternatively, maybe the best approach is to propose a split where I get the most valuable items I can, and leave the rest. Since potatoes are highest, take both. Then take as many papers or clocks as possible. Let's see:} \\
\\ 
\texttt{If I take both potatoes (2 each =4), 2 papers (1 each=2), and 1 clock (1). Total 7. The other player would get 2 papers and 1 clock} \\
            }
        }
    }
    & & \\ \\

    \theutterance \stepcounter{utterance}  
    & & \multicolumn{4}{p{0.6\linewidth}}{
        \cellcolor[rgb]{0.9,0.9,0.9}{
            \makecell[{{p{\linewidth}}}]{
                \texttt{\tiny{[GM$\rangle$P2]}}
                \texttt{<reasoning>} \\
\texttt{Let me try to figure out the best way to divide these items. The items are 4 papers, 2 clocks, and 2 potatoes. My values are paper:1, clock:1, potato:2. The other player has their own values, which I don't know. The goal is to make a proposal that gives me the highest possible score if our proposals are complementary.} \\
\\ 
\texttt{First, I need to maximize my own value. Potatoes are the most valuable to me, so I should aim to get as many as possible. There are 2 potatoes. If I take both, that's 2*2=4 points. Then, maybe take some papers or clocks. Let's see: If I take both potatoes (4), 2 papers (2), and 1 clock (1). Total 7. The other player would get 2 papers and 1 clock. Their value depends on their own values. If they value papers and clocks more than potatoes, maybe they can get more. But if they value potatoes more, they might get less. But since I don't know, maybe taking both potatoes is risky because the other player might not get anything if they also want potatoes. But if I leave one potato, maybe they take it, and we can split the rest.} \\
\\ 
\texttt{Alternatively, if I take 1 potato (2), then maybe take more papers or clocks. Let's say 1 potato, 3 papers, 1 clock. That's 2 + 3 +1 =6. The other player gets 1 potato, 1 paper, 1 clock. Their value could be higher if they value potatoes more. But I'm not sure.} \\
\\ 
\texttt{Wait, but the other player is also trying to maximize their own score. So if I take both potatoes, they might not get any, which could lead to their proposal not complementing mine if they also want potatoes. But if I leave one potato, maybe they take it, and we can split the rest.} \\
\\ 
\texttt{Alternatively, maybe the best approach is to propose a split where I get the most valuable items I can, and leave the rest. Since potatoes are highest, take both. Then take as many papers or clocks as possible. Let's see:} \\
\\ 
\texttt{If I take both potatoes (2 each =4), 2 papers (1 each=2), and 1 clock (1). Total 7. The other player would get 2 papers and 1 clock} \\
            }
        }
    }
    & \\ \\

    \theutterance \stepcounter{utterance}  
    & & & \multicolumn{4}{p{0.6\linewidth}}{
        \cellcolor[rgb]{0.9,0.9,0.9}{
            \makecell[{{p{\linewidth}}}]{
                \texttt{\tiny{[GM$\langle$P2]}}
                \texttt{<reasoning>} \\
\texttt{Let me try to work through this step by step. The goal is to maximize my score by proposing a split of the items such that my proposal and the other player's are complementary. Since I don't know their values, I need to make an educated guess based on my own preferences.} \\
\\ 
\texttt{First, let's list out the items and my values:} \\
\\ 
\texttt{{-} 4 papers, each worth 1 to me.} \\
\texttt{{-} 2 clocks, each worth 1 to me.} \\
\texttt{{-} 2 potatoes, each worth 2 to me.} \\
\\ 
\texttt{Total value if I take everything: 4*1 + 2*1 + 2*2 = 4 + 2 + 4 = 10. But obviously, the other player won't let me have everything. So I need to find a balance.} \\
\\ 
\texttt{Potatoes are the most valuable to me. There are only 2, so if I take both, that's 4 points. Then I can take some papers or clocks. Let's see: If I take both potatoes (4), 2 papers (2), and 1 clock (1), that's 7 points. The other player would get the remaining 2 papers and 1 clock. Their value depends on their preferences, but if they value papers and clocks higher than potatoes, they might get more. But if they value potatoes, they might get less. Since I don't know, maybe taking both potatoes is a good start.} \\
\\ 
\texttt{Alternatively, if I take 1 potato (2), then maybe take more papers or clocks. Let's say 1 potato, 3 papers, 1 clock. That's 2 + 3 +1 =6. The other player gets 1 potato, 1 paper, 1 clock. Their value could be higher if they value potatoes more. But I'm not sure.} \\
\\ 
\texttt{Wait, but the other player is also trying to maximize their own score. So if I take both potatoes, they might not get any, which could lead to their proposal not complementing mine if they also want potatoes. But if I leave one potato, maybe they take it, and we can split the rest.} \\
\\ 
\texttt{Alternatively, maybe the best approach is to propose a split where I get the most valuable items I can, and leave the rest. Since potatoes are highest, take both. Then take as many papers or clocks as possible. Let's see:} \\
\\ 
\texttt{If I take both potatoes (2 each =4), 2 papers} \\
            }
        }
    }
     \\ \\

    \theutterance \stepcounter{utterance}  
    & & & \multicolumn{2}{p{0.3\linewidth}}{
        \cellcolor[rgb]{0.9,0.9,0.9}{
            \makecell[{{p{\linewidth}}}]{
                \texttt{\tiny{[GM$|$GM]}}
                \texttt{continue} \\
            }
        }
    }
    & & \\ \\

    \theutterance \stepcounter{utterance}  
    & & & \multicolumn{2}{p{0.3\linewidth}}{
        \cellcolor[rgb]{0.9,0.9,0.9}{
            \makecell[{{p{\linewidth}}}]{
                \texttt{\tiny{[GM$|$GM]}}
                \texttt{<reasoning>} \\
\texttt{Let me try to work through this step by step. The goal is to maximize my score by proposing a split of the items such that my proposal and the other player's are complementary. Since I don't know their values, I need to make an educated guess based on my own preferences.} \\
\\ 
\texttt{First, let's list out the items and my values:} \\
\\ 
\texttt{{-} 4 papers, each worth 1 to me.} \\
\texttt{{-} 2 clocks, each worth 1 to me.} \\
\texttt{{-} 2 potatoes, each worth 2 to me.} \\
\\ 
\texttt{Total value if I take everything: 4*1 + 2*1 + 2*2 = 4 + 2 + 4 = 10. But obviously, the other player won't let me have everything. So I need to find a balance.} \\
\\ 
\texttt{Potatoes are the most valuable to me. There are only 2, so if I take both, that's 4 points. Then I can take some papers or clocks. Let's see: If I take both potatoes (4), 2 papers (2), and 1 clock (1), that's 7 points. The other player would get the remaining 2 papers and 1 clock. Their value depends on their preferences, but if they value papers and clocks higher than potatoes, they might get more. But if they value potatoes, they might get less. Since I don't know, maybe taking both potatoes is a good start.} \\
\\ 
\texttt{Alternatively, if I take 1 potato (2), then maybe take more papers or clocks. Let's say 1 potato, 3 papers, 1 clock. That's 2 + 3 +1 =6. The other player gets 1 potato, 1 paper, 1 clock. Their value could be higher if they value potatoes more. But I'm not sure.} \\
\\ 
\texttt{Wait, but the other player is also trying to maximize their own score. So if I take both potatoes, they might not get any, which could lead to their proposal not complementing mine if they also want potatoes. But if I leave one potato, maybe they take it, and we can split the rest.} \\
\\ 
\texttt{Alternatively, maybe the best approach is to propose a split where I get the most valuable items I can, and leave the rest. Since potatoes are highest, take both. Then take as many papers or clocks as possible. Let's see:} \\
\\ 
\texttt{If I take both potatoes (2 each =4), 2 papers} \\
            }
        }
    }
    & & \\ \\

    \theutterance \stepcounter{utterance}  
    & & \multicolumn{4}{p{0.6\linewidth}}{
        \cellcolor[rgb]{0.9,0.9,0.9}{
            \makecell[{{p{\linewidth}}}]{
                \texttt{\tiny{[P1$\langle$GM]}}
                \texttt{<reasoning>} \\
\texttt{Let me try to work through this step by step. The goal is to maximize my score by proposing a split of the items such that my proposal and the other player's are complementary. Since I don't know their values, I need to make an educated guess based on my own preferences.} \\
\\ 
\texttt{First, let's list out the items and my values:} \\
\\ 
\texttt{{-} 4 papers, each worth 1 to me.} \\
\texttt{{-} 2 clocks, each worth 1 to me.} \\
\texttt{{-} 2 potatoes, each worth 2 to me.} \\
\\ 
\texttt{Total value if I take everything: 4*1 + 2*1 + 2*2 = 4 + 2 + 4 = 10. But obviously, the other player won't let me have everything. So I need to find a balance.} \\
\\ 
\texttt{Potatoes are the most valuable to me. There are only 2, so if I take both, that's 4 points. Then I can take some papers or clocks. Let's see: If I take both potatoes (4), 2 papers (2), and 1 clock (1), that's 7 points. The other player would get the remaining 2 papers and 1 clock. Their value depends on their preferences, but if they value papers and clocks higher than potatoes, they might get more. But if they value potatoes, they might get less. Since I don't know, maybe taking both potatoes is a good start.} \\
\\ 
\texttt{Alternatively, if I take 1 potato (2), then maybe take more papers or clocks. Let's say 1 potato, 3 papers, 1 clock. That's 2 + 3 +1 =6. The other player gets 1 potato, 1 paper, 1 clock. Their value could be higher if they value potatoes more. But I'm not sure.} \\
\\ 
\texttt{Wait, but the other player is also trying to maximize their own score. So if I take both potatoes, they might not get any, which could lead to their proposal not complementing mine if they also want potatoes. But if I leave one potato, maybe they take it, and we can split the rest.} \\
\\ 
\texttt{Alternatively, maybe the best approach is to propose a split where I get the most valuable items I can, and leave the rest. Since potatoes are highest, take both. Then take as many papers or clocks as possible. Let's see:} \\
\\ 
\texttt{If I take both potatoes (2 each =4), 2 papers} \\
            }
        }
    }
    & \\ \\

    \theutterance \stepcounter{utterance}  
    & \multicolumn{4}{p{0.6\linewidth}}{
        \cellcolor[rgb]{0.9,0.9,0.9}{
            \makecell[{{p{\linewidth}}}]{
                \texttt{\tiny{[P1$\rangle$GM]}}
                \texttt{<reasoning>} \\
\texttt{Let me try to figure out the best way to divide these items. The items are 4 papers, 2 clocks, and 2 potatoes. My values are paper:1, clock:1, potato:2. The other player has their own values, which I don't know. The goal is to make a proposal that gives me the highest possible score if our proposals are complementary.} \\
\\ 
\texttt{First, I need to maximize my own value. Potatoes are the most valuable to me. There are only 2 potatoes. If I take both, that's 2*2=4 points. Then I can take some papers or clocks. Let's see: If I take both potatoes (4), 2 papers (2), and 1 clock (1). Total 7. The other player would get 2 papers and 1 clock. Their value depends on their own values. If they value papers and clocks more than potatoes, maybe they can get more. But if they value potatoes more, they might get less. Since I don't know, maybe taking both potatoes is risky because the other player might not get anything if they also want potatoes. But if I leave one potato, maybe they take it, and we can split the rest.} \\
\\ 
\texttt{Alternatively, if I take 1 potato (2), then maybe take more papers or clocks. Let's say 1 potato, 3 papers, 1 clock. That's 2 + 3 +1 =6. The other player gets 1 potato, 1 paper, 1 clock. Their value could be higher if they value potatoes more. But I'm not sure.} \\
\\ 
\texttt{Wait, but the other player is also trying to maximize their own score. So if I take both potatoes, they might not get any, which could lead to their proposal not complementing mine if they also want potatoes. But if I leave one potato, maybe they take it, and we can split the rest.} \\
\\ 
\texttt{Alternatively, maybe the best approach is to propose a split where I get the most valuable items I can, and leave the rest. Since potatoes are highest, take both. Then take as many papers or clocks as possible. Let's see:} \\
\\ 
\texttt{If I take both potatoes (2 each =4), 2 papers (1 each=2), and 1 clock (1). Total 7. The other player would get 2 papers and 1 clock. Their value depends on their preferences, but if they} \\
            }
        }
    }
    & & \\ \\

    \theutterance \stepcounter{utterance}  
    & & & \multicolumn{2}{p{0.3\linewidth}}{
        \cellcolor[rgb]{0.9,0.9,0.9}{
            \makecell[{{p{\linewidth}}}]{
                \texttt{\tiny{[GM$|$GM]}}
                \texttt{continue} \\
            }
        }
    }
    & & \\ \\

    \theutterance \stepcounter{utterance}  
    & & & \multicolumn{2}{p{0.3\linewidth}}{
        \cellcolor[rgb]{0.9,0.9,0.9}{
            \makecell[{{p{\linewidth}}}]{
                \texttt{\tiny{[GM$|$GM]}}
                \texttt{<reasoning>} \\
\texttt{Let me try to figure out the best way to divide these items. The items are 4 papers, 2 clocks, and 2 potatoes. My values are paper:1, clock:1, potato:2. The other player has their own values, which I don't know. The goal is to make a proposal that gives me the highest possible score if our proposals are complementary.} \\
\\ 
\texttt{First, I need to maximize my own value. Potatoes are the most valuable to me. There are only 2 potatoes. If I take both, that's 2*2=4 points. Then I can take some papers or clocks. Let's see: If I take both potatoes (4), 2 papers (2), and 1 clock (1). Total 7. The other player would get 2 papers and 1 clock. Their value depends on their own values. If they value papers and clocks more than potatoes, maybe they can get more. But if they value potatoes more, they might get less. Since I don't know, maybe taking both potatoes is risky because the other player might not get anything if they also want potatoes. But if I leave one potato, maybe they take it, and we can split the rest.} \\
\\ 
\texttt{Alternatively, if I take 1 potato (2), then maybe take more papers or clocks. Let's say 1 potato, 3 papers, 1 clock. That's 2 + 3 +1 =6. The other player gets 1 potato, 1 paper, 1 clock. Their value could be higher if they value potatoes more. But I'm not sure.} \\
\\ 
\texttt{Wait, but the other player is also trying to maximize their own score. So if I take both potatoes, they might not get any, which could lead to their proposal not complementing mine if they also want potatoes. But if I leave one potato, maybe they take it, and we can split the rest.} \\
\\ 
\texttt{Alternatively, maybe the best approach is to propose a split where I get the most valuable items I can, and leave the rest. Since potatoes are highest, take both. Then take as many papers or clocks as possible. Let's see:} \\
\\ 
\texttt{If I take both potatoes (2 each =4), 2 papers (1 each=2), and 1 clock (1). Total 7. The other player would get 2 papers and 1 clock. Their value depends on their preferences, but if they} \\
            }
        }
    }
    & & \\ \\

    \theutterance \stepcounter{utterance}  
    & & \multicolumn{4}{p{0.6\linewidth}}{
        \cellcolor[rgb]{0.9,0.9,0.9}{
            \makecell[{{p{\linewidth}}}]{
                \texttt{\tiny{[GM$\rangle$P2]}}
                \texttt{<reasoning>} \\
\texttt{Let me try to figure out the best way to divide these items. The items are 4 papers, 2 clocks, and 2 potatoes. My values are paper:1, clock:1, potato:2. The other player has their own values, which I don't know. The goal is to make a proposal that gives me the highest possible score if our proposals are complementary.} \\
\\ 
\texttt{First, I need to maximize my own value. Potatoes are the most valuable to me. There are only 2 potatoes. If I take both, that's 2*2=4 points. Then I can take some papers or clocks. Let's see: If I take both potatoes (4), 2 papers (2), and 1 clock (1). Total 7. The other player would get 2 papers and 1 clock. Their value depends on their own values. If they value papers and clocks more than potatoes, maybe they can get more. But if they value potatoes more, they might get less. Since I don't know, maybe taking both potatoes is risky because the other player might not get anything if they also want potatoes. But if I leave one potato, maybe they take it, and we can split the rest.} \\
\\ 
\texttt{Alternatively, if I take 1 potato (2), then maybe take more papers or clocks. Let's say 1 potato, 3 papers, 1 clock. That's 2 + 3 +1 =6. The other player gets 1 potato, 1 paper, 1 clock. Their value could be higher if they value potatoes more. But I'm not sure.} \\
\\ 
\texttt{Wait, but the other player is also trying to maximize their own score. So if I take both potatoes, they might not get any, which could lead to their proposal not complementing mine if they also want potatoes. But if I leave one potato, maybe they take it, and we can split the rest.} \\
\\ 
\texttt{Alternatively, maybe the best approach is to propose a split where I get the most valuable items I can, and leave the rest. Since potatoes are highest, take both. Then take as many papers or clocks as possible. Let's see:} \\
\\ 
\texttt{If I take both potatoes (2 each =4), 2 papers (1 each=2), and 1 clock (1). Total 7. The other player would get 2 papers and 1 clock. Their value depends on their preferences, but if they} \\
            }
        }
    }
    & \\ \\

    \theutterance \stepcounter{utterance}  
    & & & \multicolumn{4}{p{0.6\linewidth}}{
        \cellcolor[rgb]{0.9,0.9,0.9}{
            \makecell[{{p{\linewidth}}}]{
                \texttt{\tiny{[GM$\langle$P2]}}
                \texttt{<reasoning>} \\
\texttt{Let me try to work through this step by step. The goal is to maximize my score by proposing a split of the items such that my proposal and the other player's are complementary. Since I don't know their values, I need to make an educated guess based on my own preferences.} \\
\\ 
\texttt{First, let's list out the items and my values:} \\
\\ 
\texttt{{-} 4 papers, each worth 1 to me.} \\
\texttt{{-} 2 clocks, each worth 1 to me.} \\
\texttt{{-} 2 potatoes, each worth 2 to me.} \\
\\ 
\texttt{Total value if I take everything: 4*1 + 2*1 + 2*2 = 4 + 2 + 4 = 10. But obviously, the other player won't let me have everything. So I need to find a balance.} \\
\\ 
\texttt{Potatoes are the most valuable to me. There are only 2, so if I take both, that's 4 points. Then I can take some papers or clocks. Let's see: If I take both potatoes (4), 2 papers (2), and 1 clock (1). Total 7. The other player would get the remaining 2 papers and 1 clock. Their value depends on their preferences, but if they value papers and clocks higher than potatoes, they might get more. But if they value potatoes, they might get less. Since I don't know, maybe taking both potatoes is a good start.} \\
\\ 
\texttt{Alternatively, if I take 1 potato (2), then maybe take more papers or clocks. Let's say 1 potato, 3 papers, 1 clock. That's 2 + 3 +1 =6. The other player gets 1 potato, 1 paper, 1 clock. Their value could be higher if they value potatoes more. But I'm not sure.} \\
\\ 
\texttt{Wait, but the other player is also trying to maximize their own score. So if I take both potatoes, they might not get any, which could lead to their proposal not complementing mine if they also want potatoes. But if I leave one potato, maybe they take it, and we can split the rest.} \\
\\ 
\texttt{Alternatively, maybe the best approach is to propose a split where I get the most valuable items I can, and leave the rest. Since potatoes are highest, take both. Then take as many papers or clocks as possible. Let's see:} \\
\\ 
\texttt{If I take both potatoes (2 each =4), 2 papers (1} \\
            }
        }
    }
     \\ \\

    \theutterance \stepcounter{utterance}  
    & & & \multicolumn{2}{p{0.3\linewidth}}{
        \cellcolor[rgb]{0.9,0.9,0.9}{
            \makecell[{{p{\linewidth}}}]{
                \texttt{\tiny{[GM$|$GM]}}
                \texttt{continue} \\
            }
        }
    }
    & & \\ \\

    \theutterance \stepcounter{utterance}  
    & & & \multicolumn{2}{p{0.3\linewidth}}{
        \cellcolor[rgb]{0.9,0.9,0.9}{
            \makecell[{{p{\linewidth}}}]{
                \texttt{\tiny{[GM$|$GM]}}
                \texttt{<reasoning>} \\
\texttt{Let me try to work through this step by step. The goal is to maximize my score by proposing a split of the items such that my proposal and the other player's are complementary. Since I don't know their values, I need to make an educated guess based on my own preferences.} \\
\\ 
\texttt{First, let's list out the items and my values:} \\
\\ 
\texttt{{-} 4 papers, each worth 1 to me.} \\
\texttt{{-} 2 clocks, each worth 1 to me.} \\
\texttt{{-} 2 potatoes, each worth 2 to me.} \\
\\ 
\texttt{Total value if I take everything: 4*1 + 2*1 + 2*2 = 4 + 2 + 4 = 10. But obviously, the other player won't let me have everything. So I need to find a balance.} \\
\\ 
\texttt{Potatoes are the most valuable to me. There are only 2, so if I take both, that's 4 points. Then I can take some papers or clocks. Let's see: If I take both potatoes (4), 2 papers (2), and 1 clock (1). Total 7. The other player would get the remaining 2 papers and 1 clock. Their value depends on their preferences, but if they value papers and clocks higher than potatoes, they might get more. But if they value potatoes, they might get less. Since I don't know, maybe taking both potatoes is a good start.} \\
\\ 
\texttt{Alternatively, if I take 1 potato (2), then maybe take more papers or clocks. Let's say 1 potato, 3 papers, 1 clock. That's 2 + 3 +1 =6. The other player gets 1 potato, 1 paper, 1 clock. Their value could be higher if they value potatoes more. But I'm not sure.} \\
\\ 
\texttt{Wait, but the other player is also trying to maximize their own score. So if I take both potatoes, they might not get any, which could lead to their proposal not complementing mine if they also want potatoes. But if I leave one potato, maybe they take it, and we can split the rest.} \\
\\ 
\texttt{Alternatively, maybe the best approach is to propose a split where I get the most valuable items I can, and leave the rest. Since potatoes are highest, take both. Then take as many papers or clocks as possible. Let's see:} \\
\\ 
\texttt{If I take both potatoes (2 each =4), 2 papers (1} \\
            }
        }
    }
    & & \\ \\

    \theutterance \stepcounter{utterance}  
    & & \multicolumn{4}{p{0.6\linewidth}}{
        \cellcolor[rgb]{0.9,0.9,0.9}{
            \makecell[{{p{\linewidth}}}]{
                \texttt{\tiny{[P1$\langle$GM]}}
                \texttt{<reasoning>} \\
\texttt{Let me try to work through this step by step. The goal is to maximize my score by proposing a split of the items such that my proposal and the other player's are complementary. Since I don't know their values, I need to make an educated guess based on my own preferences.} \\
\\ 
\texttt{First, let's list out the items and my values:} \\
\\ 
\texttt{{-} 4 papers, each worth 1 to me.} \\
\texttt{{-} 2 clocks, each worth 1 to me.} \\
\texttt{{-} 2 potatoes, each worth 2 to me.} \\
\\ 
\texttt{Total value if I take everything: 4*1 + 2*1 + 2*2 = 4 + 2 + 4 = 10. But obviously, the other player won't let me have everything. So I need to find a balance.} \\
\\ 
\texttt{Potatoes are the most valuable to me. There are only 2, so if I take both, that's 4 points. Then I can take some papers or clocks. Let's see: If I take both potatoes (4), 2 papers (2), and 1 clock (1). Total 7. The other player would get the remaining 2 papers and 1 clock. Their value depends on their preferences, but if they value papers and clocks higher than potatoes, they might get more. But if they value potatoes, they might get less. Since I don't know, maybe taking both potatoes is a good start.} \\
\\ 
\texttt{Alternatively, if I take 1 potato (2), then maybe take more papers or clocks. Let's say 1 potato, 3 papers, 1 clock. That's 2 + 3 +1 =6. The other player gets 1 potato, 1 paper, 1 clock. Their value could be higher if they value potatoes more. But I'm not sure.} \\
\\ 
\texttt{Wait, but the other player is also trying to maximize their own score. So if I take both potatoes, they might not get any, which could lead to their proposal not complementing mine if they also want potatoes. But if I leave one potato, maybe they take it, and we can split the rest.} \\
\\ 
\texttt{Alternatively, maybe the best approach is to propose a split where I get the most valuable items I can, and leave the rest. Since potatoes are highest, take both. Then take as many papers or clocks as possible. Let's see:} \\
\\ 
\texttt{If I take both potatoes (2 each =4), 2 papers (1} \\
            }
        }
    }
    & \\ \\

    \theutterance \stepcounter{utterance}  
    & \multicolumn{4}{p{0.6\linewidth}}{
        \cellcolor[rgb]{0.9,0.9,0.9}{
            \makecell[{{p{\linewidth}}}]{
                \texttt{\tiny{[P1$\rangle$GM]}}
                \texttt{<reasoning>} \\
\texttt{Let me try to figure out the best way to divide these items. The items are 4 papers, 2 clocks, and 2 potatoes. My values are paper:1, clock:1, potato:2. The other player has their own values, which I don't know. The goal is to make a proposal that gives me the highest possible score if our proposals are complementary.} \\
\\ 
\texttt{First, I need to maximize my own value. Potatoes are the most valuable to me. There are only 2 potatoes. If I take both, that's 2*2=4 points. Then I can take some papers or clocks. Let's see: If I take both potatoes (4), 2 papers (2), and 1 clock (1). Total 7. The other player would get 2 papers and 1 clock. Their value depends on their preferences, but if they value papers and clocks higher than potatoes, maybe they can get more. But if they value potatoes, they might get less. Since I don't know, maybe taking both potatoes is risky because the other player might not get anything if they also want potatoes. But if I leave one potato, maybe they take it, and we can split the rest.} \\
\\ 
\texttt{Alternatively, if I take 1 potato (2), then maybe take more papers or clocks. Let's say 1 potato, 3 papers, 1 clock. That's 2 + 3 +1 =6. The other player gets 1 potato, 1 paper, 1 clock. Their value could be higher if they value potatoes more. But I'm not sure.} \\
\\ 
\texttt{Wait, but the other player is also trying to maximize their own score. So if I take both potatoes, they might not get any, which could lead to their proposal not complementing mine if they also want potatoes. But if I leave one potato, maybe they take it, and we can split the rest.} \\
\\ 
\texttt{Alternatively, maybe the best approach is to propose a split where I get the most valuable items I can, and leave the rest. Since potatoes are highest, take both. Then take as many papers or clocks as possible. Let's see:} \\
\\ 
\texttt{If I take both potatoes (2 each =4), 2 papers (1 each=2), and 1 clock (1). Total 7. The other player would get 2 papers and 1 clock. Their value depends on their preferences, but if they value} \\
            }
        }
    }
    & & \\ \\

    \theutterance \stepcounter{utterance}  
    & & & \multicolumn{2}{p{0.3\linewidth}}{
        \cellcolor[rgb]{0.9,0.9,0.9}{
            \makecell[{{p{\linewidth}}}]{
                \texttt{\tiny{[GM$|$GM]}}
                \texttt{continue} \\
            }
        }
    }
    & & \\ \\

    \theutterance \stepcounter{utterance}  
    & & & \multicolumn{2}{p{0.3\linewidth}}{
        \cellcolor[rgb]{0.9,0.9,0.9}{
            \makecell[{{p{\linewidth}}}]{
                \texttt{\tiny{[GM$|$GM]}}
                \texttt{<reasoning>} \\
\texttt{Let me try to figure out the best way to divide these items. The items are 4 papers, 2 clocks, and 2 potatoes. My values are paper:1, clock:1, potato:2. The other player has their own values, which I don't know. The goal is to make a proposal that gives me the highest possible score if our proposals are complementary.} \\
\\ 
\texttt{First, I need to maximize my own value. Potatoes are the most valuable to me. There are only 2 potatoes. If I take both, that's 2*2=4 points. Then I can take some papers or clocks. Let's see: If I take both potatoes (4), 2 papers (2), and 1 clock (1). Total 7. The other player would get 2 papers and 1 clock. Their value depends on their preferences, but if they value papers and clocks higher than potatoes, maybe they can get more. But if they value potatoes, they might get less. Since I don't know, maybe taking both potatoes is risky because the other player might not get anything if they also want potatoes. But if I leave one potato, maybe they take it, and we can split the rest.} \\
\\ 
\texttt{Alternatively, if I take 1 potato (2), then maybe take more papers or clocks. Let's say 1 potato, 3 papers, 1 clock. That's 2 + 3 +1 =6. The other player gets 1 potato, 1 paper, 1 clock. Their value could be higher if they value potatoes more. But I'm not sure.} \\
\\ 
\texttt{Wait, but the other player is also trying to maximize their own score. So if I take both potatoes, they might not get any, which could lead to their proposal not complementing mine if they also want potatoes. But if I leave one potato, maybe they take it, and we can split the rest.} \\
\\ 
\texttt{Alternatively, maybe the best approach is to propose a split where I get the most valuable items I can, and leave the rest. Since potatoes are highest, take both. Then take as many papers or clocks as possible. Let's see:} \\
\\ 
\texttt{If I take both potatoes (2 each =4), 2 papers (1 each=2), and 1 clock (1). Total 7. The other player would get 2 papers and 1 clock. Their value depends on their preferences, but if they value} \\
            }
        }
    }
    & & \\ \\

    \theutterance \stepcounter{utterance}  
    & & \multicolumn{4}{p{0.6\linewidth}}{
        \cellcolor[rgb]{0.9,0.9,0.9}{
            \makecell[{{p{\linewidth}}}]{
                \texttt{\tiny{[GM$\rangle$P2]}}
                \texttt{<reasoning>} \\
\texttt{Let me try to figure out the best way to divide these items. The items are 4 papers, 2 clocks, and 2 potatoes. My values are paper:1, clock:1, potato:2. The other player has their own values, which I don't know. The goal is to make a proposal that gives me the highest possible score if our proposals are complementary.} \\
\\ 
\texttt{First, I need to maximize my own value. Potatoes are the most valuable to me. There are only 2 potatoes. If I take both, that's 2*2=4 points. Then I can take some papers or clocks. Let's see: If I take both potatoes (4), 2 papers (2), and 1 clock (1). Total 7. The other player would get 2 papers and 1 clock. Their value depends on their preferences, but if they value papers and clocks higher than potatoes, maybe they can get more. But if they value potatoes, they might get less. Since I don't know, maybe taking both potatoes is risky because the other player might not get anything if they also want potatoes. But if I leave one potato, maybe they take it, and we can split the rest.} \\
\\ 
\texttt{Alternatively, if I take 1 potato (2), then maybe take more papers or clocks. Let's say 1 potato, 3 papers, 1 clock. That's 2 + 3 +1 =6. The other player gets 1 potato, 1 paper, 1 clock. Their value could be higher if they value potatoes more. But I'm not sure.} \\
\\ 
\texttt{Wait, but the other player is also trying to maximize their own score. So if I take both potatoes, they might not get any, which could lead to their proposal not complementing mine if they also want potatoes. But if I leave one potato, maybe they take it, and we can split the rest.} \\
\\ 
\texttt{Alternatively, maybe the best approach is to propose a split where I get the most valuable items I can, and leave the rest. Since potatoes are highest, take both. Then take as many papers or clocks as possible. Let's see:} \\
\\ 
\texttt{If I take both potatoes (2 each =4), 2 papers (1 each=2), and 1 clock (1). Total 7. The other player would get 2 papers and 1 clock. Their value depends on their preferences, but if they value} \\
            }
        }
    }
    & \\ \\

    \theutterance \stepcounter{utterance}  
    & & & \multicolumn{4}{p{0.6\linewidth}}{
        \cellcolor[rgb]{0.9,0.9,0.9}{
            \makecell[{{p{\linewidth}}}]{
                \texttt{\tiny{[GM$\langle$P2]}}
                \texttt{<reasoning>} \\
\texttt{Let me try to work through this step by step. The goal is to maximize my score by proposing a split of the items such that my proposal and the other player's are complementary. Since I don't know their values, I need to make an educated guess based on my own preferences.} \\
\\ 
\texttt{First, let's list out the items and my values:} \\
\\ 
\texttt{{-} 4 papers, each worth 1 to me.} \\
\texttt{{-} 2 clocks, each worth 1 to me.} \\
\texttt{{-} 2 potatoes, each worth 2 to me.} \\
\\ 
\texttt{Total value if I take everything: 4*1 + 2*1 + 2*2 = 4 + 2 + 4 = 10. But obviously, the other player won't let me have everything. So I need to find a balance.} \\
\\ 
\texttt{Potatoes are the most valuable to me. There are only 2 potatoes. If I take both, that's 4 points. Then I can take some papers or clocks. Let's see: If I take both potatoes (4), 2 papers (2), and 1 clock (1). Total 7. The other player would get the remaining 2 papers and 1 clock. Their value depends on their preferences, but if they value papers and clocks higher than potatoes, they might get more. But if they value potatoes, they might get less. Since I don't know, maybe taking both potatoes is a good start.} \\
\\ 
\texttt{Alternatively, if I take 1 potato (2), then maybe take more papers or clocks. Let's say 1 potato, 3 papers, 1 clock. That's 2 + 3 +1 =6. The other player gets 1 potato, 1 paper, 1 clock. Their value could be higher if they value potatoes more. But I'm not sure.} \\
\\ 
\texttt{Wait, but the other player is also trying to maximize their own score. So if I take both potatoes, they might not get any, which could lead to their proposal not complementing mine if they also want potatoes. But if I leave one potato, maybe they take it, and we can split the rest.} \\
\\ 
\texttt{Alternatively, maybe the best approach is to propose a split where I get the most valuable items I can, and leave the rest. Since potatoes are highest, take both. Then take as many papers or clocks as possible. Let's see:} \\
\\ 
\texttt{If I take both potatoes (2 each =4), 2 papers (1} \\
            }
        }
    }
     \\ \\

    \theutterance \stepcounter{utterance}  
    & & & \multicolumn{2}{p{0.3\linewidth}}{
        \cellcolor[rgb]{0.9,0.9,0.9}{
            \makecell[{{p{\linewidth}}}]{
                \texttt{\tiny{[GM$|$GM]}}
                \texttt{continue} \\
            }
        }
    }
    & & \\ \\

    \theutterance \stepcounter{utterance}  
    & & & \multicolumn{2}{p{0.3\linewidth}}{
        \cellcolor[rgb]{0.9,0.9,0.9}{
            \makecell[{{p{\linewidth}}}]{
                \texttt{\tiny{[GM$|$GM]}}
                \texttt{<reasoning>} \\
\texttt{Let me try to work through this step by step. The goal is to maximize my score by proposing a split of the items such that my proposal and the other player's are complementary. Since I don't know their values, I need to make an educated guess based on my own preferences.} \\
\\ 
\texttt{First, let's list out the items and my values:} \\
\\ 
\texttt{{-} 4 papers, each worth 1 to me.} \\
\texttt{{-} 2 clocks, each worth 1 to me.} \\
\texttt{{-} 2 potatoes, each worth 2 to me.} \\
\\ 
\texttt{Total value if I take everything: 4*1 + 2*1 + 2*2 = 4 + 2 + 4 = 10. But obviously, the other player won't let me have everything. So I need to find a balance.} \\
\\ 
\texttt{Potatoes are the most valuable to me. There are only 2 potatoes. If I take both, that's 4 points. Then I can take some papers or clocks. Let's see: If I take both potatoes (4), 2 papers (2), and 1 clock (1). Total 7. The other player would get the remaining 2 papers and 1 clock. Their value depends on their preferences, but if they value papers and clocks higher than potatoes, they might get more. But if they value potatoes, they might get less. Since I don't know, maybe taking both potatoes is a good start.} \\
\\ 
\texttt{Alternatively, if I take 1 potato (2), then maybe take more papers or clocks. Let's say 1 potato, 3 papers, 1 clock. That's 2 + 3 +1 =6. The other player gets 1 potato, 1 paper, 1 clock. Their value could be higher if they value potatoes more. But I'm not sure.} \\
\\ 
\texttt{Wait, but the other player is also trying to maximize their own score. So if I take both potatoes, they might not get any, which could lead to their proposal not complementing mine if they also want potatoes. But if I leave one potato, maybe they take it, and we can split the rest.} \\
\\ 
\texttt{Alternatively, maybe the best approach is to propose a split where I get the most valuable items I can, and leave the rest. Since potatoes are highest, take both. Then take as many papers or clocks as possible. Let's see:} \\
\\ 
\texttt{If I take both potatoes (2 each =4), 2 papers (1} \\
            }
        }
    }
    & & \\ \\

    \theutterance \stepcounter{utterance}  
    & & \multicolumn{4}{p{0.6\linewidth}}{
        \cellcolor[rgb]{0.9,0.9,0.9}{
            \makecell[{{p{\linewidth}}}]{
                \texttt{\tiny{[P1$\langle$GM]}}
                \texttt{<reasoning>} \\
\texttt{Let me try to work through this step by step. The goal is to maximize my score by proposing a split of the items such that my proposal and the other player's are complementary. Since I don't know their values, I need to make an educated guess based on my own preferences.} \\
\\ 
\texttt{First, let's list out the items and my values:} \\
\\ 
\texttt{{-} 4 papers, each worth 1 to me.} \\
\texttt{{-} 2 clocks, each worth 1 to me.} \\
\texttt{{-} 2 potatoes, each worth 2 to me.} \\
\\ 
\texttt{Total value if I take everything: 4*1 + 2*1 + 2*2 = 4 + 2 + 4 = 10. But obviously, the other player won't let me have everything. So I need to find a balance.} \\
\\ 
\texttt{Potatoes are the most valuable to me. There are only 2 potatoes. If I take both, that's 4 points. Then I can take some papers or clocks. Let's see: If I take both potatoes (4), 2 papers (2), and 1 clock (1). Total 7. The other player would get the remaining 2 papers and 1 clock. Their value depends on their preferences, but if they value papers and clocks higher than potatoes, they might get more. But if they value potatoes, they might get less. Since I don't know, maybe taking both potatoes is a good start.} \\
\\ 
\texttt{Alternatively, if I take 1 potato (2), then maybe take more papers or clocks. Let's say 1 potato, 3 papers, 1 clock. That's 2 + 3 +1 =6. The other player gets 1 potato, 1 paper, 1 clock. Their value could be higher if they value potatoes more. But I'm not sure.} \\
\\ 
\texttt{Wait, but the other player is also trying to maximize their own score. So if I take both potatoes, they might not get any, which could lead to their proposal not complementing mine if they also want potatoes. But if I leave one potato, maybe they take it, and we can split the rest.} \\
\\ 
\texttt{Alternatively, maybe the best approach is to propose a split where I get the most valuable items I can, and leave the rest. Since potatoes are highest, take both. Then take as many papers or clocks as possible. Let's see:} \\
\\ 
\texttt{If I take both potatoes (2 each =4), 2 papers (1} \\
\\ 
\\ 
\texttt{The time is up. It is now your turn to submit a secret proposal.} \\
            }
        }
    }
    & \\ \\

    \theutterance \stepcounter{utterance}  
    & \multicolumn{4}{p{0.6\linewidth}}{
        \cellcolor[rgb]{0.9,0.9,0.9}{
            \makecell[{{p{\linewidth}}}]{
                \texttt{\tiny{[P1$\rangle$GM]}}
                \texttt{<reasoning>} \\
\texttt{Let me try to figure out the best proposal. The items are 4 papers, 2 clocks, 2 potatoes. My values are paper:1, clock:1, potato:2. The other player's values are unknown, but I need to maximize my score.} \\
\\ 
\texttt{Potatoes are the most valuable to me. Taking both gives me 4 points. Then, I can take papers or clocks. Let's see: If I take both potatoes, 2 papers, and 1 clock, that's 4 + 2 + 1 =7. The other player would get the remaining 2 papers and 1 clock. Their value might be lower if they value potatoes more. But if they value papers/clocks more, they might get more. But since I don't know, maybe this is a safe split.} \\
\\ 
\texttt{Alternatively, taking 1 potato (2), 3 papers (3), and 1 clock (1) gives me 6. The other player gets 1 potato, 1 paper, 1 clock. But if they value potatoes, they might prefer that. But I'm not sure.} \\
\\ 
\texttt{Wait, but the other player is also trying to maximize their own score. If I take both potatoes, they might not get any, which could make their proposal not complementary. But if I leave one potato, maybe they take it. Let's compare:} \\
\\ 
\texttt{Option 1: Both potatoes, 2 papers, 1 clock. My score: 4+2+1=7. Their score: 2 papers +1 clock =3. Total items: 2p+2p+2p+1c+1c= Wait, no. Wait, total items are 4p, 2c, 2p. If I take 2p, 1c, 2p (wait, no, potatoes are 2 total). Wait, correction: I take both potatoes (2), 2 papers, 1 clock. That's 2 potatoes, 2 papers, 1 clock. Remaining items: 2 papers, 1 clock. So their proposal must include those. Their value would be 2*1 +1*1=3. My value is 2*2 +2*1 +1*1=4+2+1=7. That's good for me.} \\
\\ 
\texttt{Option 2: Take 1 potato, 3 papers, 1 clock.} \\
            }
        }
    }
    & & \\ \\

    \theutterance \stepcounter{utterance}  
    & & & \multicolumn{2}{p{0.3\linewidth}}{
        \cellcolor[rgb]{0.9,0.9,0.9}{
            \makecell[{{p{\linewidth}}}]{
                \texttt{\tiny{[GM$|$GM]}}
                \texttt{continue} \\
            }
        }
    }
    & & \\ \\

    \theutterance \stepcounter{utterance}  
    & & & \multicolumn{2}{p{0.3\linewidth}}{
        \cellcolor[rgb]{0.9,0.9,0.9}{
            \makecell[{{p{\linewidth}}}]{
                \texttt{\tiny{[GM$|$GM]}}
                \texttt{player was instructed to make a proposal} \\
            }
        }
    }
    & & \\ \\

\end{supertabular}
}

\end{document}
